\documentclass{article}
\usepackage{circuitikzgit}
\usetikzlibrary{tikzmark,fpu}

\tikzset{ov/.pic = {
% reference anchor is -center
    \draw coordinate(-center)
    (-center) +(-1.2,-1) rectangle +(1.2,1)
    (-center) ++(-1.2,0.8) coordinate (-in 1)
    (-center) ++(-1.2,-0.8) coordinate (-in 2)
    (-center) ++(1.2,0.8) coordinate (-out 1)
    (-center) ++(1.2,-0.8) coordinate (-out 2)
    (-center) ++(0,1) coordinate (-up)
    (-in 1) -- ++(0.5,0) coordinate(-tmp)
    to[leD*, diodes/scale=0.6, led arrows from cathode, name=-led]
    (-tmp|- -in 2) -- (-in 2)
    (-out 1) -- ++(-0.5,0) coordinate(-tmp)
    to[pD*, diodes/scale=0.4, mirror, name=-pd1] ++(0,-0.5)
    edge[densely dashed] ++(0,-0.533) ++(0,-0.566)
    to[pD*, diodes/scale=0.4,mirror, name=-pd2] (-tmp|- -out 2) -- (-out 2)
% leave the position of the path at the center
    (-center);
    }
}
\begin{document}

Circuitikz: \pgfcircversion{} (\pgfcircversiondate)\par
Ti\emph{k}Z: \pgfversion{} (\pgfversiondate)\par\bigskip


\begin{tikzpicture}
    \draw (0,0) pic(one) {ov};
    \node [above] at (one-up) {O1};
    \draw[color=blue] (one-out 1) -- ++(1,0)
        % pic[pic anchor=(-in 1)](two) {ov};
        pic[tikzmark round=3, rotate=45, transform shape, pic anchor=(-in 1)](two) {ov};
    \node [above] at (two-up) {O2};
\end{tikzpicture}
\end{document}
